problem:  
Let \((M^{2n+1}, \eta, \xi, \Phi, g)\) be a compact *irregular* Sasakian manifold, i.e. the Reeb vector field \(\xi\) generates an *irregular* Reeb foliation (the orbits are not all closed). Let \(\Delta_b\) denote the sub-Laplacian associated with the Tanaka–Webster connection and the canonical contact volume form \(\eta \wedge (d\eta)^n\). Define the **transverse first eigenvalue** of the sub-Laplacian by  
\[
\lambda_1^T = \inf\left\{ \frac{\int_M |\nabla_b f|^2 \, d\mu}{\int_M f^2\, d\mu} : f \in C^\infty(M), \int_M f\, d\mu = 0 \right\},
\]
where \(\nabla_b\) denotes the horizontal gradient.

Suppose it is known that for any *quasi-regular* Sasakian structure on \(M\), the corresponding basic Kähler quotient \(Z=M/\mathcal{F}_\xi\) satisfies a **Lichnerowicz-type estimate** \(\lambda_1^T \geq c(n,\rho)\) depending on the lower bound \(\rho\) of the transverse Ricci curvature.  

**Question:**  
Does there exist an *irregular* Sasakian structure \((\eta,\xi,\Phi,g)\) with the same underlying smooth manifold and with
\[
\lambda_1^T = c(n,\rho),
\]
i.e., achieving the same spectral extremum as some quasi-regular structure, but with a *dense Reeb flow*?  
More generally, can such a spectral rigidity occur in the closure of the Calabi cone over \(M\), regarded as a limit of quasi-regular structures, or does spectral extremality intrinsically enforce quasi-regularity of the Reeb foliation?

Why is it a "good" problem:  
This problem explicitly ventures beyond the classical spectral rigidity territory addressed by CR and pseudo-Hermitian Obata–Lichnerowicz theorems, which almost all assume compact *regular* or *quasi-regular* Sasakian manifolds where the Reeb foliation has smooth Kähler quotient. It asks whether spectral extremality of the sub-Riemannian Laplacian—a global analytic condition—forces geometric regularity of the Reeb dynamics. Thus it links three domains rarely intertwined in a concrete question:  
(1) **sub-Riemannian spectral geometry** (through \(\lambda_1^T\));  
(2) **contact and Sasakian foliation theory** (through irregular Reeb dynamics); and  
(3) **transverse Kähler geometry** (through the role of quasi-regular approximations).

The question is simple to state yet deeply unresolved: there is currently no known mechanism connecting spectral data of the horizontal Laplacian to the arithmetic or dynamical nature (rational or irrational) of the Reeb flow. Resolving even special cases would require fundamentally new analytic insights binding dynamical invariants, CR analysis, and Sasakian structure moduli—qualifying it as a “good” high-level research problem.